\section{Product cells and logarithmic packing}
\label{b:product-section}

There is another way to exploit opposite eigenvalues.  Expansion in one
coordinate is accompanied by contraction in the other, so their product
provides a useful bound for an intercept patch.  A cell with a long, thin
arm in each coordinate direction can absorb these changes at arbitrarily
small relative area cost.  We prove this using open cylinder bases and
interior packing.  This method requires the boundary of the parameter
set to have area zero.

\begin{theorem}[Approximation by product cells]
\label{b:product}
Let $h,M>0$, let $D\subset\mathbb R^2$ be compact with
$|\partial D|=0$, and let $V$ be $C^1$ on an open neighborhood of $D$.
Suppose that $DV(p)$ has eigenvalues $\lambda_p,-\lambda_p$, where
$0<M\lambda_p<1$ for every $p\in D$.  For every $e>0$, the set
\[
 \mathcal S_h(D,V,M)=
 \{(x,p_1+xV_1(p),h(p_2+xV_2(p))):p\in D,\ |x|\le M\}
\]
has a finite cylinder cover with bounded open measurable perpendicular
bases whose total area is at most
\[
 h\int_D J_h(V(p))\,dp+e,
 \qquad J_h(w)=(1+w_1^2+h^2w_2^2)^{-1/2}.
\]
\end{theorem}

In the notation of Section~\ref{b:introduction},
$\mathcal S_h(D,V,M)=F_h(E(D,V,M))$.

We first calculate the cell geometry.  The same calculation also
explains its logarithmic parametrization, which will be used below with
a different stopping rule for the packing.

\begin{lemma}[Cell area and dilation]
\label{b:cell-geometry}
For $A>0$ and $0<B\le A^2$, the set
\[
 H(A,B)=\{(u_1,u_2):|u_1|,|u_2|\le A,\ |u_1u_2|\le B\}
\]
has area $4B(1+\log(A^2/B))$.  Strict inequalities give the same area.
For $0<\tau<1$ and $\gamma\ge0$, define
\[
 \begin{split}
 Q_\tau&=\{s:|s_1|,|s_2|\le1,\ |s_1s_2|\le\tau\},\\
 P_{\tau,\gamma}&=\{s:|s_1|,|s_2|<3,
                         \ |s_1s_2|<(1+\gamma)\tau\}.
 \end{split}
\]
If $(1+\gamma)\tau<1$, then
\begin{equation}
 \frac{|P_{\tau,\gamma}|}{|Q_\tau|}
 =R(\tau,\gamma)
 :=\frac{(1+\gamma)[1+\log(9/((1+\gamma)\tau))]}
          {1+\log(1/\tau)}.
 \label{b:cell-ratio}
\end{equation}
For $m>1$, put
\[
 S_m=H(m,1),\qquad
 S'_{m,\gamma}=\{u:|u_1|,|u_2|<3m,
                            \ |u_1u_2|<1+\gamma\}.
\]
With $\tau=m^{-2}$, one has the exact identities
\begin{equation}
 S_m=mQ_\tau,\qquad S'_{m,\gamma}=mP_{\tau,\gamma}.
 \label{b:cell-dilation}
\end{equation}
\end{lemma}

\begin{proof}
The area in the first quadrant is
\[
 \int_0^{B/A}A\,du_1+
 \int_{B/A}^{A}\frac{B}{u_1}\,du_1
 =B+B\log(A^2/B).
\]
The boundaries have area zero, so replacing weak inequalities by strict
ones leaves the result unchanged.  Formula~\eqref{b:cell-ratio} follows
by applying this identity to $H(1,\tau)$ and $H(3,(1+\gamma)\tau)$.
Substitution $u=ms$ proves~\eqref{b:cell-dilation}.
\end{proof}

Given $\eta>0$, there are two useful orders for obtaining
$R(\tau,\gamma)<1+\eta$.  One can first choose
$0<\gamma<\eta$ and then take $\tau$ small, because
$R(\tau,\gamma)\to1+\gamma$ as $\tau\downarrow0$.
Alternatively, first take $m$ large so that
$R(m^{-2},0)<1+\eta$, and then take $\gamma>0$ small.
In either case the spatial scale is chosen only after the shape and
padding.  These choices make the open intercept patch almost as cheap
as the closed parameter cell.

\begin{lemma}[One cell and its cylinder]
\label{b:product-local}
Assume the hypotheses of Theorem~\ref{b:product}, with $D$ nonempty.
For each $p\in D$, let $L_p$ have unit eigenvectors for
$\lambda_p,-\lambda_p$ as its columns.  Fix $\eta>0$,
$0<\tau<1$, and $\gamma>0$ such that $(1+\gamma)\tau<1$ and
$R(\tau,\gamma)\le1+\eta$.  There is $\rho_0>0$ such that, whenever
$0<\rho<\rho_0$ and
\[
 G=p+\rho L_pQ_\tau\subset D,
\]
the open cylinder with axis direction $(1,V_1(p),hV_2(p))$ and
parameter intercept patch
\[
 U_{p,\rho}=p+\rho L_pP_{\tau,\gamma}
\]
covers all segments labeled by $G$ in $\mathcal S_h(D,V,M)$.
Its perpendicular base has area at most
\begin{equation}
 h(1+\eta)^2\int_G J_h(V(q))\,dq.
 \label{b:product-local-cost}
\end{equation}
Here a parameter intercept $u$ denotes the physical point
$(0,u_1,hu_2)$.
\end{lemma}

\begin{proof}
There are constants $\beta>0$ and $N<\infty$ such that
$|\det L_p|\ge\beta$ and $\|L_p^{-1}\|\le N$ for all $p\in D$,
independently of the choices of unit eigenvectors.  Otherwise a convergent
subsequence of centers and unit eigenvectors would give parallel
limiting eigenvectors.  The continuous positive function
$\lambda_p=\sqrt{-\det DV(p)}$ has a positive minimum on $D$,
so those limiting vectors would belong to distinct eigenspaces of the
same matrix.  This is impossible.  Also $\|L_p\|\le\sqrt2$.

Choose $a>0$ so that the compact set
$D_a=\{q:\operatorname{dist}(q,D)\le a\}$ lies in the domain of $V$.
Let $\omega$ be a modulus of continuity for $DV$ on $D_a$.
If $q=p+\rho L_ps$, $s\in Q_\tau$, and $2\rho<a$, the entire
segment from $p$ to $q$ lies in $D_a$, and Taylor's integral formula gives
\[
 |V(q)-V(p)-DV(p)(q-p)|\le2\rho\,\omega(2\rho).
\]
The line through a point labeled by $q,x$, parallel to the chosen
axis, meets the intercept plane at parameter
\begin{equation}
 q+x(V(q)-V(p))
 =p+\rho L_p\big(( (1+x\lambda_p)s_1,
                         (1-x\lambda_p)s_2)+E\big),
 \qquad |E|\le2MN\omega(2\rho).
 \label{b:product-taylor}
\end{equation}
For the linear part $u$ inside the parentheses, the spectral hypothesis
and $|x|\le M$ give $|u_i|\le2$ and
\[
 |u_1u_2|=|(1-x^2\lambda_p^2)s_1s_2|\le\tau.
\]
Write $\epsilon_\rho=2MN\omega(2\rho)$.  When
$\epsilon_\rho<\min\{1/2,\gamma\tau/5\}$, one has
\[
 |u_i+E_i|<3,\qquad
 |(u_1+E_1)(u_2+E_2)|
 \le\tau+4\epsilon_\rho+\epsilon_\rho^2
 <(1+\gamma)\tau.
\]
Thus all the intercepts lie in $U_{p,\rho}$, uniformly in the center
and the segment parameter.

Set $j(p)=J_h(V(p))$.  By Proposition~\ref{b:projection}, the
perpendicular base obtained from this intercept patch has area
\[
 h j(p)|U_{p,\rho}|
 =h j(p)R(\tau,\gamma)|G|
 \le h(1+\eta)j(p)|G|.
\]
The function $j$ is positive and uniformly continuous on $D$.
After reducing $\rho_0$, it satisfies
$j(p)\le(1+\eta)j(q)$ for $q\in G$; this proves
\eqref{b:product-local-cost}.  The projection from the physical
intercept plane onto the perpendicular plane is an invertible linear
map, since the axis has nonzero first coordinate.  The base is
therefore bounded and open, as asserted.  No continuous choice of
eigenvectors is needed.
\end{proof}

We have constructed a cylinder for each sufficiently small cell, with
cost controlled by the integral on that cell.  To sum these estimates,
we place disjoint cells in the interior and then cover the omitted
labels separately.

\begin{proof}[Proof of Theorem~\ref{b:product}]
For $D=\varnothing$ use no cylinders.  Otherwise choose $\eta>0$ so
that
\[
 h[(1+\eta)^2-1]\int_D J_h(V)<e/2,
\]
and choose $\gamma$, then $\tau$, then $\rho_0$ as above.
In any bounded open remainder $O\subset\operatorname{int}D$ of
positive area, take a sufficiently fine square grid of side $\ell$
whose closed squares contained in $O$ account for at least $|O|/2$.
Such grids exist at arbitrarily fine scales: every point of $O$ belongs
to a contained grid square once its diameter is smaller than the
point's distance to the complement, and bounded convergence gives the
area assertion.  There are only finitely many retained squares.

At the center $p$ of each square put $\rho=\ell/8<\rho_0$.
Since $|L_ps|\le |s_1|+|s_2|\le2$ for $s\in Q_\tau$, the cell
$G=p+\rho L_pQ_\tau$ lies in the disk of radius $\ell/4$ about $p$,
hence strictly inside the square.  Its area is at least
$\beta|Q_\tau|\ell^2/64$.  The cells are disjoint.  Remove their
finite compact union and repeat in the open remainder.  Each round
removes at least the fraction
\[
 c=\beta|Q_\tau|/128>0
\]
of the remaining area.  Consequently the area remaining after $n$
rounds is at most $(1-c)^n|D|$.  There are at most countably many
selected cells, and they cover the interior except for a null set.
Disjointness and~\eqref{b:product-local-cost} bound the total cost of
their open cylinders by $h(1+\eta)^2\int_D J_h(V)$.

Let $R$ be the labels in $D$ not contained in these cells.
It includes any uncovered boundary points, and $|R|=0$ because
$|\partial D|=0$.  Cover $R$ by countably many open squares with
sum of areas less than $\sigma$, where $\sigma>0$ will be chosen
below, and with uniformly small side lengths.  The latter restriction
is available by subdividing an outer-measure cover and slightly
enlarging the resulting squares within the area allowance.
Discard squares not meeting $R$.  In each remaining square of side
$s$, choose $p\in R$.  If $K$ bounds $\|DV\|$ on the compact
neighborhood $D_a$ from the preceding proof, then for every $q\in R$
in that square and $|x|\le M$,
\[
 |q+x(V(q)-V(p))-p|\le(1+MK)\sqrt2\,s,
\]
provided $\sqrt2s<a$.  Indeed the segment joining $p$ to $q$
stays in $D_a$, so the derivative bound applies even when the domain
of $V$ is not convex.  An open intercept disk about $p$ of radius
$2(1+MK)\sqrt2s$ covers all these intercepts.  Its associated
perpendicular base costs at most
\[
 C_0s^2,\qquad C_0=8\pi h(1+MK)^2,
\]
because $J_h(V(p))\le1$.  Choose $C_0\sigma<e/2$.

All the resulting cylinders are open and together cover the compact
set $\mathcal S_h(D,V,M)$.  A finite subcover has no larger total
cost, proving the claimed bound.  If $|D|=0$, use only the residual
disk construction; its cost can be arbitrarily small.
\end{proof}

\subsection{Logarithmic cells and finite stopping}
\label{b:logarithmic-section}

The dilation in Lemma~\ref{b:cell-geometry} permits a second choice of
shape parameters and a finite stopping rule.  We state the resulting
formulation explicitly because it will be useful when a smooth field
is multiplied by a small tilt parameter.

\begin{corollary}[Approximation by logarithmic cells]
\label{b:logarithmic}
Let $X,k>0$, let $U\subset\mathbb R^2$ be compact with
$|\partial U|=0$, and let $v$ be $C^2$ on an open rectangle containing
$U$.  Suppose that, for every $p\in U$,
\[
 \operatorname{tr}Dv(p)=0,\qquad \det Dv(p)<0,\qquad
 kX\sqrt{-\det Dv(p)}<1.
\]
In orthonormal coordinates $(x,p)\in\mathbb R\times\mathbb R^2$,
the set
\[
 \{(x,p+xkv(p)):p\in U,\ |x|\le X\}
\]
has, for every $e>0$, a finite cylinder cover with bounded open
measurable perpendicular bases and total base area at most
\[
 \int_U(1+k^2|v(p)|^2)^{-1/2}\,dp+e.
\]
\end{corollary}

\begin{proof}
Set $V=kv$, $h=1$ and $M=X$ in Theorem~\ref{b:product}.
The eigenvalues of $DV(p)$ are $\pm k\sqrt{-\det Dv(p)}$,
so that theorem gives the stated conclusion.  We record the
logarithmic parametrization and its finite stopping rule.

Put $j(p)=(1+k^2|v(p)|^2)^{-1/2}$ and choose $\eta>0$ with
$[(1+\eta)^2-1]\int_U j<e/2$.  As after
Lemma~\ref{b:cell-geometry}, choose $m>1$ and then $\gamma>0$
with $R(m^{-2},\gamma)<1+\eta$ and $(1+\gamma)m^{-2}<1$.
For the unit-eigenvector matrix $L_p$, the cells
\[
 G=p+tL_pS_m,\qquad U_{p,t}=p+tL_pS'_{m,\gamma}
\]
are exactly the product cell and its open patch under
$\tau=m^{-2}$ and $\rho=mt$.  Lemma~\ref{b:product-local}
therefore supplies an open cylinder with axis parallel to $(1,kv(p))$
and cost at most $(1+\eta)^2\int_G j$ whenever $G\subset U$
and $t$ is sufficiently small, uniformly in $p$.

For fixed $m$, Taylor's theorem bounds the error in the label-plane
intercept by $O_m(t^2)$, uniformly in $p$ and $|x|\le X$.
After applying $t^{-1}L_p^{-1}$, this is $O_m(t)$.  The
corresponding linear image of $S_m$ has coordinate bounds $2m$
and absolute product bound $1$, so
this error fits the strict margins $3m$ and $1+\gamma$ once $t$
is small.  These estimates use a fixed compact rectangle containing
$U$ in its interior and contained in the domain of $v$.

Use the preceding square packing with $t=\ell/(8m)$.
Each round removes at least the fraction
$\beta|S_m|/(128m^2)=\beta|Q_{m^{-2}}|/128>0$
of the open remainder.  Stop after finitely many rounds when its
area is below $\alpha$.  Because $|\partial U|=0$, the entire
residual label set $R\subset U$ then has area below $\alpha$;
the finitely many cell cylinders cost at most
$(1+\eta)^2\int_U j$.

Cover $R$ by countably many open squares whose areas sum to less
than $2\alpha$.
If $K_v$ bounds $Dv$ on the fixed rectangle, the residual-disk
construction in the preceding proof costs less than
$16\pi(1+XkK_v)^2\alpha$.  Choose $\alpha>0$ to make this less
than $e/2$.  The cylinders form an open cover of the compact swept
set, so a finite subcover proves the bound.  For $|U|=0$, use only
the residual disks; for $U=\varnothing$, use no cylinders.
\end{proof}

The two endings preserve different finite constructions: logarithmic
packing stops with finitely many cells and a small-area residual,
whereas product packing first exhausts the interior up to a null set.
Both finish with residual disks and compactness.  The resulting bases
are bounded and open, and the cell bases need not be convex.  The
field, physical scale and axial interval are fixed before any shape,
mesh or finite-subcover choices are made.
